\subsection{X-ray lenses}


In visible-light optics the refractive index of a material is greater than
unity ($n>1$), meaning light slows down inside the medium and bends toward the
surface normal at an interface.  For X-rays the situation is reversed: the
refractive index is slightly \emph{less} than unity and is written as
%
\begin{equation}
  n = 1 - \delta + i\beta,
  \label{eq:refindex}
\end{equation}
%
where $\delta$ is the \textbf{refractive index decrement} (a small positive
number, typically $10^{-5}$--$10^{-7}$) and $\beta$ is the \textbf{absorption
index}.  Because $\delta$ is so small, refraction effects are extremely
weak---a single interface deflects an X-ray beam by only micro-radians.  This
is why refractive lenses for X-rays were considered impractical for so long and
why focusing by mirrors (total external reflection), zone plates, and bent
crystals was preferred instead.

Two important consequences follow directly from $n < 1$:
\begin{itemize}
  \item A \textbf{concave} lens acts as a converging lens for X-rays (the
        opposite of visible light), because the beam bends away from the normal
        at an interface going from a denser to a less-dense medium.
  \item There is a \textbf{critical angle} for total external reflection: for
        incidence angles shallower than
        $\theta_c \approx \sqrt{2\delta}$, X-rays are totally reflected rather
        than transmitted.  This limits the maximum useful numerical aperture of
        a refractive lens.
\end{itemize}

\paragraph{The Compound Refractive Lens (CRL)}
Because $\delta$ is so small, a single concave lens has an impractically long
focal length.  For aluminium at $14\,\mathrm{keV}$,
$\delta \approx 2.8\times10^{-6}$, giving a single-lens focal length
%
\begin{equation}
  F = \frac{R}{2\delta} \approx 54\,\mathrm{m}
\end{equation}
%
for a hole radius $R = 0.3\,\mathrm{mm}$.  The solution proposed by
Snigirev~et~al.\ \\cite{Snigirev1996} is the \textbf{compound refractive lens
(CRL)}: an array of $N$ identical concave lenses aligned along a common optical
axis.  The focal length of the compound system is
%
\begin{equation}
  F = \frac{R}{2N\delta},
  \label{eq:focal}
\end{equation}
%
so stacking 30 lenses reduces the focal length to a manageable $\sim1.8\,\mathrm{m}$.
% This works because each lens contributes an equal, additive deflection to any
% ray passing through it, in direct analogy with a stack of weak prisms.  The
% compound lens can then be treated as a single thin lens obeying the
% \textbf{Gaussian lens formula}:
% %
% \begin{equation}
%   \frac{1}{p} + \frac{1}{q} = \frac{1}{F},
%   \label{eq:gauss}
% \end{equation}
% %
% where $p$ is the source distance and $q$ is the image distance.  

The
diffraction-limited resolution of the lens is
\begin{equation}
  \sigma_f = \frac{\lambda\, q}{A},
\end{equation}
%
where $A$ is the effective aperture and $\lambda$ is the X-ray wavelength. Note that the diffraction limited size is not considered by ray tracing.  

\paragraph{Lens Surface Shapes and Aberrations}

The choice of surface profile has a strong effect on focus quality, especially
at high numerical aperture.
Spherical lenses are the easiest to manufacture but suffer from severe
\textbf{spherical aberration}: marginal rays (passing through the lens edge)
are focused at a different position from paraxial rays (passing near the
optical axis), producing a large, diffuse focal spot.  X-ray spherical lenses
are no longer used in practice.
Parabolic lenses eliminate most spherical aberration and have been the standard
choice since the late 1990s.  A parabolic profile $z^2 = 2Ry$ closely
approximates the ideal shape for most beamline configurations, and parabolic
CRLs routinely achieve micrometre-scale focus sizes. However, the parabola is
not a perfect focusing surface.

Elliptical lenses are the theoretically ideal surface for focusing a
\textbf{collimated} beam to a point.  The ellipse is a conic section obtained
as a special case of the Cartesian oval when the source is at infinity.
Ray-tracing simulations \cite{sanchezdel2012} show that for nanofocusing
applications, elliptical profiles outperform parabolic ones not in terms of the
FWHM of the central peak, but in the \textbf{intensity of the tails}
surrounding it.  Those tails reduce the peak intensity (and therefore the
gain), so at large apertures the ellipse is clearly preferred.

The \textbf{Cartesian oval} is the general ideal surface shape for
\textbf{point-to-point focusing} from a source at finite distance.
The Cartesian oval is a bi-circular quartic
curve and reduces to an ellipse (for $n_1 < n_2$, as in X-ray optics) or
hyperbola (for $n_1 > n_2$, as in visible-light glass optics) when the source
is collimated.  In high-demagnification nanofocusing geometries, the Cartesian
oval produces a clean Gaussian focal spot with no tails, while elliptical
lenses already show some tails and parabolic lenses show significant ones.
Table~\ref{tab:shapes} summarises the ideal surface choice depending on beam
and focusing geometry for X-ray optics ($n < 1$).

\begin{table}[h!]
  \centering
  \caption{Ideal refractive surface shape for different focusing conditions
           (X-rays, $n<1$).}
  \label{tab:shapes}
  \begin{tabular}{lll}
    \toprule
    \textbf{Input beam} & \textbf{Output beam} & \textbf{Ideal surface} \\
    \midrule
    Collimated          & Convergent (focused) & Ellipse                \\
    Convergent          & Collimated           & Hyperbola              \\
    Divergent           & Convergent           & Cartesian oval         \\
    \bottomrule
  \end{tabular}
\end{table}


\textbf{SHADOW} works for any x-ray lens  simulation, and proceeds as follows.
\begin{itemize}
    \item Each ray is propagated in a straight line from the source to the first lens surface.  The coordinates and direction cosines of the ray at the intersection with the surface are calculated analytically (for simple shapes like paraboloids) or numerically (for Cartesian ovals, using a pre-computed surface mesh.
    \item At every air--material interface, the Snell's law in vector form is applied. This is repeated at every interface of every lens in the array.
    \item After traversing a path length $l_i$ through the lens material, the electric vectors of the $i$-th ray are affected by absorption  $e^{-\mu_i\, l_i}$, where $\mu_i$ is the linear absorption coefficient at the ray's energy, including photoelectric absorption and Compton scattering.
\end{itemize}

SHADOW can output the full two-dimensional intensity distribution at any plane, from which one extracts the focused spot size, the effective aperture, or compute chromatic aberrations under polychromatic illumination. Obviously, the effect of upstream optics such as monochromators or mirrors can be included in the simulations. 

The main limitation of geometric ray tracing is the absence of diffraction. For very small focal spots (below $\sim100\,\mathrm{nm}$) or high-NA configurations, diffraction broadening becomes comparable to the geometric spot size and must be added separately, typically by convolving the ray-tracing point-spread function with an Airy disc of diameter $\sim\lambda/\mathrm{NA}$. The ``hybrid" method, that corrects the ray tracing of some deffractive effects is available in OASYS and can be used for that.